\documentclass[12pt]{article}
\usepackage[utf8]{inputenc}
\usepackage[T1]{fontenc}
\usepackage{hyperref}
\title{Manombai Language Resources for K-12 Curriculum Development}
\author{Compiled by Claude Sonnet 5.5 (claude-sonnet-5-5) for Andr\'as Kornai}
\date{}
\begin{document}
\maketitle

\section{Introduction}
Manombai (ISO 639-3: woo; Glottocode: mano1275) is an Austronesian language of the Aru Islands, Maluku, Indonesia. Search results give it the alternative names Wokam, Woka, Manobai and Wamar. Glottolog places it in the Central Aru subgroup of the Aru branch of Malayo-Polynesian. Speaker counts differ by source: about 9,100 (Wikipedia, 2011), 4,490 (ASJP) and 7,475 (Global Recordings Network). Very little material exists, and I found no school materials in the language.

\section{Primary Educational Resources}
I found no primers, textbooks or school programmes for Manombai in search results.

\section{Linguistic and Reference Materials}
\begin{itemize}
\item Wikipedia article on the Manombai language: \url{https://en.wikipedia.org/wiki/Manombai_language}
\item ASJP database entry, which holds a basic word list: \url{https://asjp.clld.org/languages/MANOMBAI}
\item Endangered Languages Project bibliography for Manombai: \url{https://endangeredlanguages.com/lang/4648/bibliography}
\end{itemize}

\section{Religious and Community Resources}
No scripture translation or community organisation was found in search results.

\section{Digital and Audio Resources}
\begin{itemize}
\item Global Recordings Network language page for woo (audio recordings may be listed there; I did not verify their content): \url{https://globalrecordings.net/en/language/woo}
\end{itemize}

\section{Curriculum Development Context}
The documentation I found is thin: reference entries, a word list and a bibliography. A curriculum would first need new documentation, such as an orthography, a lexicon and recorded texts, developed with speaker communities in the Aru Islands. Indonesian would be the likely bridge language.
\end{document}